%!TEX TS-program = xelatex
\documentclass[]{friggeri-cv}
\usepackage{afterpage}
\usepackage{hyperref}
\usepackage{color}
\usepackage{xcolor}
\usepackage{smartdiagram}
\usepackage{fontspec}
% if you want to add fontawesome package
% you need to compile the tex file with LuaLaTeX
% References:
%   http://texdoc.net/texmf-dist/doc/latex/fontawesome/fontawesome.pdf
%   https://www.ctan.org/tex-archive/fonts/fontawesome?lang=en
%\usepackage{fontawesome}
\usepackage{metalogo}
\usepackage{dtklogos}
\usepackage[utf8]{inputenc}
\usepackage{tikz}
\usetikzlibrary{mindmap,shadows}
\hypersetup{
    pdftitle={Florent Clarret - CV},
    pdfauthor={Florent Clarret},
    pdfsubject={CV},
    pdfkeywords={Clarret, resume, cv},
    colorlinks=false,           % no lik border color
    allbordercolors=white       % white border color for all
}
\smartdiagramset{
    bubble center node font = \footnotesize,
    bubble node font = \footnotesize,
    % specifies the minimum size of the bubble center node
    bubble center node size = 0.5cm,
    %  specifies the minimum size of the bubbles
    bubble node size = 0.5cm,
    % specifies which is the distance among the bubble center node and the other bubbles
    distance center/other bubbles = 0.3cm,
    % sets the distance from the text to the border of the bubble center node
    distance text center bubble = 0.5cm,
    % set center bubble color
    bubble center node color = pblue,
    % define the list of colors usable in the diagram
    set color list = {lightgray, materialcyan, orange, green, materialorange, materialteal, materialamber, materialindigo, materialgreen, materiallime},
    % sets the opacity at which the bubbles are shown
    bubble fill opacity = 0.6,
    % sets the opacity at which the bubble text is shown
    bubble text opacity = 0.5,
}

\addbibresource{bibliography.bib}
\RequirePackage{xcolor}
\definecolor{pblue}{HTML}{0395DE}

\begin{document}

\header{Florent}{ Clarret}
      {Ingénieur étude et développement}
      
% Fake text to add separator      
\fcolorbox{white}{gray}{\parbox{\dimexpr\textwidth-2\fboxsep-2\fboxrule}{%
.....
}}



% In the aside, each new line forces a line break
\begin{aside}
  \includegraphics[scale=0.18]{img/photo.png}
  \section{Adresse}
    10 Rue des barattières
    37130, Cinq-Mars-La-Pile
    France
    ~
\section{Infos}
    Né le 01/04/1992
    Nationalité française
    Permis B et véhicule
    ~
  \section{Contact}
   +33 (0)6 62 97 01 69
    \href{mailto:florent.clarret@gmail.com}{\textbf{florent.clarret@}\\gmail.com}
    ~
  \section{Web}
	    \href{https://www.linkedin.com/in/florent-clarret/}{florent-clarret} \includegraphics{img/linkedin.png} 
	    \href{https://github.com/FlorentClarret}{FlorentClarret} \includegraphics{img/github.png} 
    ~
  \section{Langues}
	Français
    Anglais
    ~
  \section{Développement}
    \smartdiagram[bubble diagram]{
        \textbf{Java},
        \textbf{Shell},
        \textbf{Spring}\\\textbf{Boot},
        \textbf{JPA},
        \textbf{Tomcat},
        \textbf{Active}\\\textbf{MQ},
        \textbf{Linux},
        \textbf{Oracle},
        \textbf{MySQL}
    }
    ~
  \section{Outillage}
    \smartdiagram[bubble diagram]{
        \textbf{Tools},
        \textbf{Maven},
        \textbf{Gitlab},
        \textbf{Docker},
        \textbf{CI/CD},
        \textbf{Git},
        \textbf{Dynatrace},
        \textbf{Jmeter},
        \textbf{JMH}
    }
    ~
\end{aside}
~
\\
\\
\section{Expériences professionnelles}
\begin{entrylist}
      \vspace{7pt}    
  \entry
    {Depuis 05/18}
    {Ingénieur étude et développement Java}
    {Worldline, Tours}
    {Refonte technique et évolution fonctionnelle du Back-End d'une solution d'authentificiation forte des personnes. Ajout de moyens alternatifs d'authentification par fingerprint ou le protocole FIDO en plus de l'utilisation du PIN. Mise en place d'une nouvelle API permettant une communication client/serveur plus efficace et sécurisée. Optimisation  du Back-End afin de pouvoir traiter toujours plus d'authentifications et toujours plus rapidement. Passage de Spring 3 vers Spring Boot 2, le tout s'appuyant sur une base de données de type MySQL. \\Dockerisation de l'application pour faciliter la mise en place des tests d'intégration.}
      \vspace{7pt}    
  \entry
    {Depuis 09/17}
    {Enseignant}
    {Polytech Tours}
    {Dispense de cours pour les étudiants de dernière année du cursus d'ingénieur portant sur le développement en Java, son écosystème ainsi que sur l'aspect performance (64h par an).}
      \vspace{7pt}    
    \entry
    {05/15 - 05/18}
    {Ingénieur étude et développement Java}
    {Worldline, Tours}
    {Etude, développement et déploiement en production de la nouvelle version du Back-End de la solution de paiement sur internet SIPS sur une technologie Java avec une base de données Oracle. Mise  en place des nouveaux traitements pour assurer la sauvegarde et le traitement des transactions bancaires. Ecriture et optimisation de nouveaux modules pour absorber le traffic grandissant. Communication avec des services tiers via ActiveMQ et web services.\\ Participation active à la migration en production des clients sur cette nouvelle solution. Mise en place de surveillances spécifiques des modules via ELK.\\ Maintenance de la première version de la solution lors de la phase de transition. Les traitements étaient réalisés majoritairement en SQL (Sybase, MySQL) et C ainsi qu'un ensemble de scripts Shell et awk.\\Communication principalement en anglais (Benelux, Arménie et Inde).}
      \vspace{7pt}      
    \entry
    {09/13 - 09/15}
    {Etudiant-enseignant}
    {Polytech Tours}
    {Cours personnalisés aux étudiants de première année de cycle d'ingénieur. Ils portaient sur l'algorithmique et les structures de données ainsi que l'implémentation de ceux-ci en langage C (2h par semaine).}
      \vspace{7pt}      
    \entry
    {Etés 10/11/12}
    {Agent administratif intérimaire}
    {ANGDM, Lens}
    {}
\end{entrylist}

\section{Formation}
\begin{entrylist}
  \vspace{-6pt}    
  \entry
    {2012 - 2015}
    {Diplôme d'ingénieur en informatique}
    {Polytech Tours}
    {}
  \vspace{-6pt}    
  \entry
    {2014 - 2014}
    {Diplôme d'ingénieur en informatique}
    {Politechnika Łódzka, Pologne}
    {}
  \entry
    {2010 - 2012}
    {Diplôme universitaire de technologie en informatique}
    {IUT de Lens}
    {}
\end{entrylist}

\section{Centres d'intérêts}
\begin{entrylist}
  \entry
    {Homelaber. Lecture (Fantasy \& mangas). Course à pieds}
    {}    {}    {}
\end{entrylist}

\end{document}
